% Exam Template for UMTYMP and Math Department courses
%
% Using Philip Hirschhorn's exam.cls: http://www-math.mit.edu/~psh/#ExamCls
%
% run pdflatex on a finished exam at least three times to do the grading table on front page.
%
%%%%%%%%%%%%%%%%%%%%%%%%%%%%%%%%%%%%%%%%%%%%%%%%%%%%%%%%%%%%%%%%%%%%%%%%%%%%%%%%%%%%%%%%%%%%%%

% These lines can probably stay unchanged, although you can remove the last
% two packages if you're not making pictures with tikz.
\documentclass[11pt]{exam}
\RequirePackage{amssymb, amsfonts, amsmath, latexsym, verbatim, xspace, setspace}
\RequirePackage{tikz, pgflibraryplotmarks}

% By default LaTeX uses large margins.  This doesn't work well on exams; problems
% end up in the "middle" of the page, reducing the amount of space for students
% to work on them.
\usepackage[margin=1in]{geometry}
\usepackage{enumerate}
\usepackage{amsthm}

\theoremstyle{definition}
\newtheorem{soln}{Solution}

% Here's where you edit the Class, Exam, Date, etc.
\newcommand{\class}{Math 150A Section 8 and 16}
\newcommand{\term}{Spring 2025}
\newcommand{\examnum}{Practice Exam III Solutions}
\newcommand{\examdate}{March 10, 2025}
\newcommand{\timelimit}{1 Hour 50 Minutes}

% For an exam, single spacing is most appropriate
\singlespacing
% \onehalfspacing
% \doublespacing

% For an exam, we generally want to turn off paragraph indentation
\parindent 0ex

\begin{document} 

% These commands set up the running header on the top of the exam pages
\pagestyle{head}
\firstpageheader{}{}{}
\runningheader{\class}{\examnum\ - Page \thepage\ of \numpages}{\examdate}
\runningheadrule

\begin{flushright}
\begin{tabular}{p{2.8in} r l}
\textbf{\class} & \textbf{Name (Print):} & \makebox[2in]{\hrulefill}\\
\textbf{\term} &&\\
\textbf{\examnum} & \textbf{Student ID:}&\makebox[2in]{\hrulefill}\\
\textbf{\examdate} &&\\
\textbf{Time Limit: \timelimit} % & Teaching Assistant & \makebox[2in]{\hrulefill}
\end{tabular}\\
\end{flushright}
\rule[1ex]{\textwidth}{.1pt}


This exam contains \numpages\ pages (including this cover page) and
\numquestions\ problems.  Check to see if any pages are missing.  Enter
all requested information on the top of this page, and put your initials
on the top of every page, in case the pages become separated.\\

You may \textit{not} use your books or notes on this exam.  However, you may use a single, handwritten, one-sided notesheet and a \textit{basic} calculator.\\

You are required to show your work on each problem on this exam.  The following rules apply:\\

\begin{minipage}[t]{3.7in}
\vspace{0pt}
\begin{itemize}

%\item \textbf{If you use a ``fundamental theorem'' you must indicate this} and explain
%why the theorem may be applied.

\item \textbf{Organize your work}, in a reasonably neat and coherent way, in
the space provided. Work scattered all over the page without a clear ordering will 
receive very little credit.  

\item \textbf{Mysterious or unsupported answers will not receive full
credit}.  A correct answer, unsupported by calculations, explanation,
or algebraic work will receive no credit; an incorrect answer supported
by substantially correct calculations and explanations might still receive
partial credit.  This especially applies to limit calculations.

\item If you need more space, use the back of the pages; clearly indicate when you have done this.

\item \textbf{Box Your Answer} where appropriate, in order to clearly indicate what you consider the answer to the question to be.
\end{itemize}

Do not write in the table to the right.
\end{minipage}
\hfill
\begin{minipage}[t]{2.3in}
\vspace{0pt}
%\cellwidth{3em}
\gradetablestretch{2}
\vqword{Problem}
\addpoints % required here by exam.cls, even though questions haven't started yet.	
\gradetable[v]%[pages]  % Use [pages] to have grading table by page instead of question

\end{minipage}
\newpage % End of cover page

%%%%%%%%%%%%%%%%%%%%%%%%%%%%%%%%%%%%%%%%%%%%%%%%%%%%%%%%%%%%%%%%%%%%%%%%%%%%%%%%%%%%%
%
% See http://www-math.mit.edu/~psh/#ExamCls for full documentation, but the questions
% below give an idea of how to write questions [with parts] and have the points
% tracked automatically on the cover page.
%
%
%%%%%%%%%%%%%%%%%%%%%%%%%%%%%%%%%%%%%%%%%%%%%%%%%%%%%%%%%%%%%%%%%%%%%%%%%%%%%%%%%%%%%

\begin{questions}

\addpoints

\question[10]\mbox{}
\textbf{TRUE or FALSE!}  Write 
\begin{enumerate}[(a)]
\item  $$\int x\cos(x)dx = x\sin(x) + \cos(x) + C$$
TRUE
\item  $$\sum_{k=1}^4 \frac{1}{k} = \frac{25}{12}$$
TRUE
\item  $1 + 2 + 3 + 4 + \dots + 100 = 5050$
\\
TRUE
\item  if $\int_0^1 f(x)dx = 1$ then $\int_0^1 3f(x) dx = 3$
\\
TRUE
\item  assuming $f$ and $g$ are continuous
$$\int f(x)g(x)dx = \left(\int f(x)dx\right)\left(\int g(x)dx\right)$$
FALSE
\end{enumerate}

\newpage
\question[10]\mbox{}
\begin{enumerate}[(a)]
\item  Calculate $h'(x)$ for $h(x)$ the function
$$h(x) = \int_{x^3}^x \tan(1/t) dt$$
We use the formula
$$\frac{d}{dx}\int_{u(x)}^{v(x)} f(t)dt = f'(v(x))v'(x)-f'(u(x))u'(x)$$
to get
$$h'(x) = \tan(1/x)(x)' - \tan(1/x^3)(x^3)' = \tan(1/x)-3x^2\tan(1/x^3).$$
\item  Calculate $\int_0^3 (\sqrt{9-x^2} + x + 1)dx$ by drawing a picture
We break it up into three integrals:
$$\int_0^3 (\sqrt{9-x^2} + x + 1)dx = \int_0^3 \sqrt{9-x^2}dx + \int_0^3 xdx + \int_0^3 1dx.$$
This consists of a quarter circle of radius $3$, a triangle of base $3$ and height $3$, and a rectangle of base $3$ and height $1$.  Picture not included...
The integral is then
$$\int_0^3 (\sqrt{9-x^2} + x + 1)dx  = \frac{1}{4}(9\pi) + \frac{1}{2}3^2 + 3.$$
\end{enumerate}

\newpage
\question[10]\mbox{}
Use Newton's method to complete the table below and approximate a solution of the equation accurate to four decimal places

$$x + \cos(x) = 0.$$

Newton's method says to use the update equation:
$$x_{n+1} = x_n - \frac{f(x)}{f'(x)},$$
where here $f(x) = x +\cos(x)$.
Therefore for our particular situation, the update is
$$x_{n+1} = x_n - \frac{x_n+\cos(x_n)}{1-\sin(x_n)},$$
Thus
$$x_2 = -1.000000 - \frac{-1.000000+\cos(-1.000000)}{1-\sin(-1.000000)} = -0.750364.$$
$$x_3 = -0.750364 - \frac{-0.750364+\cos(-0.750364)}{1-\sin(-0.750364)} = -0.739113.$$
$$x_4 = -0.739113 - \frac{-0.739113+\cos(-0.739113)}{1-\sin(-0.739113)} = -0.739085.$$
$$x_5 = -0.739085 - \frac{-0.739085+\cos(-0.739085)}{1-\sin(-0.739085)} = -0.739085.$$


\begin{center}
\begin{tabular}{c|r}
$x_0$  &  $0.000000$\\\hline
$x_1$  &  $-1.000000$\\\hline
$x_2$  &  $-0.750364$\\\hline
$x_3$  &  $-0.739113$\\\hline
$x_4$  &  $-0.739085$\\\hline
$x_5$  &  $-0.739085$
\end{tabular}
\end{center}

\newpage
\question[10]\mbox{}
Compute each of the following integrals.  Show all work.
Lack of work results in no credit!!
\begin{enumerate}[(a)]
\item $\int_1^2 u^3 + 3u + 1 du$
\begin{align*}
\int_1^2 u^3 + 3u + 1 du
  & = \left(\frac{1}{4}u^4 + \frac{3}{2}u^2 + u\right) |_1^2\\
  & = \left(\frac{1}{4}2^4 + \frac{3}{2}2^2 + 2\right) - \left(\frac{1}{4}1^4 + \frac{3}{2}1^2 + 1\right) = \frac{37}{4}.
\end{align*}
\item $\int (x-1)(x^2+2) dx$
\begin{align*}
\int (x-1)(x^2+2) dx
  & = \int x^3 - x^2 + 2x-2 dx\\
  & = \frac{1}{4}x^4 - \frac{1}{3}x^3 + x^2 - 2x + C
\end{align*}
\item $\int_1^8 \frac{3z^3 + \sqrt[3]{z}}{z}dz$
\begin{align*}
\int_1^8 \frac{3z^3 + \sqrt[3]{z}}{z}dz
  & = \int_1^8 3z^2 + z^{-2/3} dz\\
  & = (z^3 + 3z^{1/3})|_1^8\\
  & = (8^3 + 3\cdot 8^{1/3}) - (1^3 + 3\cdot 1^{1/3}) = 514
\end{align*}
\item $\int \csc(t)\cot(t)dt$
Since the derivative of $\csc(t)$ is $-\csc(t)\cot(t)$, we get
$$\int \csc(t)\cot(t)dt = -\csc(t) + C.$$
\end{enumerate}

\newpage
\question[10]\mbox{}

A company wishes to design a fish tank in the shape of a box with an open top, where the length of the base is twice the width.
The sides of the tank are glass, which costs $3$ dollars per square meter and the base is made out of metal which costs $5$ dollars per square meter.
If the volume of the tank needs to be $10$ cubic meters, what dimensions will minimize the cost of production.

Let the dimensions of the box be $x, y, z$, where $x$ is the width, $y$ is the length, and $z$ is the height.
Then the area of the base is $xy$, so the cost of the base is
$$C_{\text{base}} = 5xy.$$
The area of the four sides combined is $xz + yz + xz + yz$, so the cost of the sides is
$$C_{\text{sides}} = 3(xz + yz + xz + yz) = 6xz + 6yz.$$
Therefore the total cost is
$$C = 5xy + 6xz + 6yz.$$
Now we know that the length of the base is twice the width, so $y = 2x$.
Therefore our cost function is 
$$C = 10x^2 + 6xz + 12xz = 10x^2 + 18xz.$$
Finally, the volume needs to be $10$ cubic meters, so 
$$10 = xyz = 2x^2z.$$
Therefore we see that $z = 5x^{-2}$.
Thus our cost is
$$C = 10x^2 + 90x^{-1}.$$
To find the dimensions minimizing this cost, we need to find a critical point.
We take the derivative and set it equal to zero, getting
$$20x - 90x^{-2} = 0.$$
Solving this, we get $x = (9/2)^{1/3}$.  Plugging this into our equations for $y$ and $z$, we get that the maximum dimensions are
$$x = (9/2)^{1/3},\ \ \ y = 2(9/2)^{1/3},\ \ \ z = 5(2/9)^{2/3}.$$

\newpage
\question[10]\mbox{}

\begin{enumerate}[(a)]
\item  Determine an integral whose value is equal to $\lim_{n\rightarrow\infty} R_n$ where here
$$R_n = \sum_{k=1}^n \sin\left(1 + \frac{3k}{n}\right) \frac{3}{n}.$$
The Riemann sum for the integral $\int_a^b f(x)dx$ is
$$R_n = \sum_{k=1}^n f(x_k)\Delta x,\ \ \ \Delta x = \frac{b-a}{n},\ \ x_k = a + k\Delta x.$$
From this, we can identify $a=1$, $f(x) = \sin(x)$, and $\Delta x = \frac{3}{n}$.
Therefore $\frac{b-a}{n} = \frac{3}{n}$, so that $b-a=3$ and since $a =1$ this says $b=4$.
This says the integral we want is $\int_1^4 \sin(x)dx$.
\item  Compute the value of the left Riemann sum $L_3$ for $f(x) = \cos(x)$ on the interval $[-\pi/3,\pi/3]$.
The width of our rectangles is $\Delta x = 2\pi/9$.
Therefore using the left endpoints of the rectangles
$$L_3 = \cos(-\pi/3)(2\pi/9) + \cos(-\pi/9)(2\pi/9) + \cos(\pi/9)(2\pi/9).$$
\end{enumerate}

\newpage
\question[10]\mbox{}

Consider the Riemann integral

$$\int_0^2 x^2 + x dx$$

\begin{enumerate}[(a)]
\item  Determine an explicit expression for the right Riemann sum $R_n$
The widths of the rectangles are $\Delta x = \frac{2-0}{n} = 2/n$.
Therefore $x_k = 0 + k\Delta_x = 2k/n$.
The right Riemann sum is
$$R_n = \sum_{k=1}^n f(x_k)\Delta x = \sum_{k=1}^n \left(\left(\frac{2k}{n}\right)^2+ \left(\frac{2k}{n}\right)\right)\frac{2}{n}.$$
\item  Compute Riemann integral by calculating the limit $\lim_{n\rightarrow\infty} R_n$
We can rewrite the above sum as
\begin{align*}
R_n
  &= \sum_{k=1}^n \left(\left(\frac{2k}{n}\right)^2+ \left(\frac{2k}{n}\right)\right)\frac{2}{n}\\
  &= \sum_{k=1}^n \frac{8k^2}{n^3} + \frac{4k}{n^2}\\
  &= \sum_{k=1}^n \frac{8k^2}{n^3} + \sum_{k=1}^n \frac{4k}{n^2}\\
  &= \frac{8}{n^3}\sum_{k=1}^nk^2 + \frac{4}{n^2}\sum_{k=1}^n k = \frac{8}{n^3}\frac{n(n+1)(2n+1)}{6} + \frac{4}{n^2}\frac{n(n+1)}{2}.
\end{align*}
Taking the limit, we get
\begin{align*}
\lim_{n\rightarrow\infty} R_n
  &= \lim_{n\rightarrow\infty}\left(\frac{8(2n^3+3n^2+n)}{6n^3} + \frac{4n^2+4n}{2n^2}\right)\\
  &= \lim_{n\rightarrow\infty}\frac{8(2n^3+3n^2+n)}{6n^3} + \lim_{n\rightarrow\infty}\frac{4n^2+4n}{2n^2}\\
  &= \frac{16}{6} + \frac{4}{2} = \frac{28}{6}.
\end{align*}
\end{enumerate}



\end{questions}
\end{document}
